\paragraph{Richiesta (a)}
Sulla base di semplici argomenti di natura qualitativa (ovvero discutendone la risposta ai
limiti del dominio delle frequenze reali senza effettuare il calcolo della funzione di
trasferimento, come richiesto invece al punto successivo), mostrare che sia il blocco 1
che l’intero circuito svolgano effettivamente la funzione di filtro passa-basso e calcolare
per ciascuno il guadagno nel limite di continua.

\paragraph{Richiesta (b)}
Derivare analiticamente le funzioni di trasferimento $A_1(s)=V_1(s)/V_S(s)$,
$A_2(s)=V_2(s)/V_1(s)$ nonché la funzione di trasferimento complessiva $A(s)$,
individuando per ciascuna poli e zeri (si assuma che gli operazionali siano
ideali e per semplicit\'a che R1=R2=R3=R5=R=10kW, R2=R/10=1kW e C1=C2=C=100nF),
e verificare che gli andamenti osservati al punto 1a siano in accordo.

\paragraph{Risposta (b)}
Un'amplificatore invertente con opamp e impedenze ha una soluzione ben nota
\begin{equation*}
  H_1(s) = -\frac{R_2 || \frac{1}{sC_1}}{R_1} = -\frac{1}{1 + sRC}
\end{equation*}
\begin{equation*}
  H_2(s) = -\frac{R_5}{\left(R_3 || \frac{1}{sC_2}\right) + R_4} = -10\frac{1 + sRC}{11 + sRC}
\end{equation*}
\begin{equation}
  \label{eq:tfunc-compensazione-polo-zero}
  \colorbox{blu}{\boxed{H(s) = \frac{10}{11 + sRC}}}
\end{equation}

\paragraph{Richiesta (c)}
Dalle espressioni ottenute dedurre i valori attesi per i parametri dei filtri realizzati
rispettivamente dal blocco 1 e dall’intero circuito (frequenze di taglio, guadagno
massimo) e confrontarli con le misure di cui al punto 1b.

\paragraph{Richiesta (d)}
Determinare l’andamento temporale atteso per i segnali in V1 e V2, verificando che
descrivano con sufficiente accuratezza le forme d’onda osservate al punto 1c.

\paragraph{Richiesta (e)}
Calcolare i valori attesi per il tempo di discesa di V1 e di salita di V2 e confrontarli con
quanto misurato al punto 1d.

\paragraph{Richiesta (f)}
Illustrare qualitativamente l’effetto, sia nel dominio temporale che delle frequenze,
dell’applicazione del blocco 2 al segnale in V1.


